\section{继续筛查发现的关键问题：相同速率不等于相同重联}
\label{sec:remaining-findings}
本节把此次15篇新增核读组织为问题链。它们包含直接Hall/双流体研究、混合模型、PIC、电阻MHD和实验，模型身份均保留；不能因用于Hall综述就把所有发现改称Hall模拟发现。

\subsection{全局重联是否存在系统尺寸效应}
\citet{F600}在固定宏观形状的二维Harris片中发现：尺寸增加30倍，PIC的全局平均率约降低2.6倍；Hall扫描增加100倍，率约降低3.4倍。小尺寸Hall结果近似$\ell^{-1/3}$，更大尺寸下降变缓，是否最终平台仍未确定。排气张角变窄伴随率下降；PIC持续产生磁岛，而Hall后期近单X点，Hall中的波列也没有在PIC中同样出现。

这对“跨尺度时一个固定率就足够”提出直接限制，但不构成对所有局地0.1机制的简单否定。初始片厚是否随盒子一起变、比较局地还是远上游场、按物理时间还是相同磁通阶段统计，会改变被回答的问题。第一章压力亏损模型针对局地准稳态；本篇针对整个演化的磁通处理，二者需要在同一诊断下重新比较。该文的平均区间排印矛盾及质量比25使其精确指数仍应审慎使用。

\subsection{导引场有时提高率，也会改变模型之间的差距}
磁岛合并研究\citep{L241}的结果比“导引场抑制重联”更丰富。小系统中从无导引场增加到$B_g=B_0$，PIC与混合模型峰值率分别提高约51\%和53\%，而Hall MHD略降。此时动理学模型中的离子更磁化，离子层变薄，堆积场和喷流增强，Hall与动理学模型的宏观率趋于接近。再增大导引场，则率重新缓慢下降。

无导引场时，PIC离子喷流比Hall结果弱约七倍，离子非旋转对称应力是重要阻力；Hall标量压力不能表达它。导引场较强时虽然宏观喷流相近，压力与应力的各项分解仍不相同。十矩用合适闭合长度可在小系统接近PIC，较大无导引场系统却产生PIC没有的次级岛，说明闭合的可移植性仍是问题。

全局岛合并还会反弹：若片内处理磁通不够快，两个岛先接近，再分开，之后再合并。无导引场PIC平均率随岛尺寸下降较陡；$B_g=B_0$时四种模型的平均率指数接近。大尺度无导引场PIC堆积区出现镜模相容的磁峰和密度凹陷；作者的线性色散检验支持该解释，但它是否进一步\emph{导致}重联减慢尚未证明。

\subsection{次级磁岛并不总带来更快的净能量释放}
三个证据应并列而非互相替代。第一，磁岛合并的强导引场PIC会不断生成次级岛，Hall却维持开放X形，二者平均率仍相近\citep{L241}。第二，PIC质量比扫描中低质量比容易生成岛，但岛并未使全局能量转换像电阻MHD算例那样显著提高；主电子层和移动次级X点的非理想预算也不同\citep{F606}。第三，欠分辨的MHD可以产生额外岛和假平台率\citep{F788}。

因此值得研究的新问题是：在什么驱动、寿命和尺度区间，岛的生成真正控制磁通通量？不能仅用“有岛”推断“岛介导的率”，也不能用一个无岛的已分辨算例否定其他初值下的撕裂。质量比趋于真实值、离子FLR、导引场、噪声和开放边界仍需联动检验。

\subsection{电子单独重联可以由大尺度重联和剪切流二次生成}
F655的2026年预印本\citep{F655}给出两条清楚的时序候选机制。其一，大尺度离子耦合重联先形成岛链，喷流使岛互相挤压，小于离子耦合尺度的新片再出现电子喷流。其二，湍流先生成电子速度剪切，随后出现与电子Kelvin--Helmholtz不稳定性相容的涡链，磁场被卷曲后形成小重联点。电子单独重联因此不必全靠直接在离子附近尺度注入能量。

这些图像有研究价值，但“太阳风中普遍发生”和“主导加热”尚超出本算例能证明的范围。系统仍为二维周期域，注入尺度约$50d_i$；电子尺度仅约一格。算法找到516个强电流结构，其中69\%薄于$d_i$且短于$5d_i$，并未逐个证实都是重联。部分$D_e$为负，所有片的平均量很小；不能把强局地正峰相加便当作净不可逆耗散。

结合解析速率研究\citep{F741}还得到另一个容易忽略的现象：无离子喷流并不意味着离子不反馈。周期小系统的离子反压仍可通过面内电场限制电子平台速度，使远上游离子归一率为$O(1)$；电子峰值、平台速度、局地电子归一率是三个不同量。

\subsection{Full-F、任意波长极化与多X点不同步}
GREENY及其后续论文\citep{E473,L314}显示，保留任意波长极化后，某些算例比长波近似发展更快并出现更细的电势和电流结构；全密度模型还保留了小振幅模型忽略的乘积项。初始电流由电子还是离子承担会明显改变过程，后续研究中离子载流启动通常慢约一个数量级。拉长盒子后可有多个X点在不同阶段同时活动，不能用一条单X点曲线代表整个域。

非正规线性算子提供暂态放大的可能，非线性阶段又伴随谱展宽；两者不能直接画等号。暖离子会改变态的形貌，但主岛数不随温比稳定单调增加。若只看摘要中的“新模型更快”或“非正规性导致爆发”，容易漏掉正文对范数、排序、初值和参数标签的限制。第一章给出了可以自行检查的能量范数和极化推导。

\subsection{Hall发电机既能协助，也能抵消磁松弛}
扩展MHD的RFP研究\citep{F798}不是孤立二维X点问题，而是三维多模撕裂和平均电流松弛。首次准单螺旋态崩解时，Hall与MHD发电机同向展平电流；后续多螺旋事件中，Hall项有时抵消部分MHD项。主模减弱、邻近模增强改变了相关电动势的符号；净松弛仍持续。

相应Lorentz力先驱动动量，Reynolds响应随后抵消一部分。后续事件中核心速度与电流的改变量同向，更接近MST集合观测；只看首次事件会得到相反判断。这个现象说明“Hall必然加速”过于笼统，也说明只验证总磁能释放相近不足以验证离子流和动量输运。

\subsection{多组分带来新的Hall图样和电场抵消}
冷离子MMS事件\citep{L227}表明，不同离子群可在同一分离线具有不同磁化程度；冷离子仍随电场漂移，其对流项可与Hall项同量级。负尘埃PIC\citep{L283}进一步出现内外极性相反的嵌套四极Hall场，外侧结构随着缓慢尘埃响应而建立。后者是三物种结果，不能列成标准电子--离子两流体已解释的现象；尘埃质量比、电荷和密度只取少数组合，也没有对应实验确认。
